\section{Introduction and Task Definition}
\label{sec:intro}

\paragraph{The task as a black box.}
DiCo-NLI \citep{dico-nli-2027-task} gives a system an ordered pair of short phrases, a premise $p$ and a hypothesis $h$, and asks for one of four labels, the set $\mathcal{L}$: \EQ{} (the phrases mean the same), \FW{} ($p$ entails $h$), \BW{} ($h$ entails $p$) or \NO{} (none of these; \cref{tab:task}). It is single-label, four-way classification over ordered pairs. Every reversible source pair appears twice, once in each order, and the gold labels of the two orders are tied by a fixed reversal operator: \EQ{} stays \EQ{}, and \FW{} and \BW{} swap. The organisers report three scores side by side: weighted F1 over all instances, \SoftCons{} (do the system's two answers on a pair obey the reversal operator) and \HardCons{} (are both answers correct); \cref{sec:evaluation} defines them. What a system observes is the two phrase strings and their language codes; the instance and pair identifiers are bookkeeping.

\begin{table}[t]
\centering
\small
\begin{tabular}{@{}p{0.27\columnwidth}p{0.65\columnwidth}@{}}
\toprule
Input & ordered phrase pair $(p, h)$, English, about three words each \\
Output & one label per ordered pair \\
Label space & \EQ, \FW, \BW, \NO \\
Formulation & four-way classification with a cross-instance constraint between $(p,h)$ and $(h,p)$ \\
Official scores & weighted F1, \SoftCons, \HardCons \\
Our data & track 1 (English): train 3,042, dev 660 instances; dev serves as our test set \\
\bottomrule
\end{tabular}
\caption{DiCo-NLI track 1 at a glance.}
\label{tab:task}
\end{table}

\paragraph{The task as a white box.}
Whether one phrase entails another depends on whether its meaning is contained in the other's: in natural logic, adding a restrictive modifier narrows a phrase and yields forward entailment \citep{maccartney-manning-2009-extended}. The relevant cues are therefore which phrase carries extra material, whether the heads match, and what quantifiers, numerals and prepositions do to scope. The same structure invites shortcuts: if the more specific phrase is usually the longer one, length and word containment can predict direction without any semantics, and \cref{sec:eda} shows they do.

Three further confounds matter. The instance identifier leaks the label, since only reversible pairs have a \texttt{\_\_flipped} twin, so no model input may include it. A fifth of the development phrases also occur in training. The source of each pair (image captions or news headlines) is not part of the release.

\paragraph{Why the task is hard.}
The phrases are about three words long and carry no sentence context. Separating equivalence from entailment needs lexical knowledge (``a chocolate bar'' and ``a candy bar''), negatives are often a single substituted word inside an identical frame (``off Mexico'' and ``on Mexico''), and consistency is a property of two predictions that a classifier makes independently.

\paragraph{Question and findings.}
We ask whether systems reach directional consistency by representing the relation between the two phrases or by exploiting surface asymmetries, and how much of a fine-tuned encoder's three scores a cue-only model recovers. We test three hypotheses. H1: a cue-only model reaches a large share of the encoder's \SoftCons{}, because its cues are symmetric or change sign exactly under reversal. H2: the encoder's advantage concentrates where the cues are uninformative, including the boundary between equivalence and entailment. H3: accuracy and consistency trade off in models trained on single ordered instances. H1 and H2 hold; H3 holds inside our classical baselines and fails for the encoder.
